\section{Every mixed-prime spectral direction and its actual pulse receiver}


This paper independently derives the spectrum and every finite pulse
eigenspace proposed in Sections 6--8 of the supplied continuation
\emph{Full-support reconstruction: the missing cohomology, the actual
prime degree, and complete mixed-sign calculations}.  The calculation
retains the original support diagonal, all $s!$ tensor summands, the
original pulse amplitudes, and the actual map from a compression back
to the full Hilbert space.  It also derives every higher tensor
evaluation on the same pulse tests.  The finite positivity estimate
used in the last section is the proved companion result \cite{FinitePrimePulse};
no positivity or boundedness of a full Weil operator is assumed.

The analytic source is Alain Connes, \emph{The Riemann Hypothesis:
Past, Present and a Letter Through Time}, arXiv:2602.04022v1,
original author source \path{rhready.tex}, subsection
\texttt{sectweilpos}, equations \texttt{wp} and \texttt{arch}
(source lines 841--868).  Connes attributes the criterion there to
Andr\'e Weil.  The exact programme antecedent is the pinned
\href{https://github.com/KokunoYumeto/zeta-function-research-reader/blob/3c78cfbd9e194d35117b33277c421c6564ce3dbf/workbenches/splitzero-tandem/continuations/20260923-full-prime-support/033/MIXED_PRIME_ENERGY_AND_BOUNDARY.tex#L209}{edition 033, PME6--10},
with the original three-pulse calculation
\href{https://github.com/KokunoYumeto/zeta-function-research-reader/blob/3c78cfbd9e194d35117b33277c421c6564ce3dbf/workbenches/splitzero-tandem/continuations/20260923-full-prime-support/033/MIXED_PRIME_ENERGY_AND_BOUNDARY.tex#L648}{FWP2--23}
and the shared-Archimedean support trace
\href{https://github.com/KokunoYumeto/zeta-function-research-reader/blob/3c78cfbd9e194d35117b33277c421c6564ce3dbf/workbenches/splitzero-tandem/continuations/20260923-full-prime-support/033/MIXED_PRIME_ENERGY_AND_BOUNDARY.tex#L1005}{FWJ1--8}.
These antecedents and the supplied Sections 6--8 are identified
separately from the new derivation below.  No prime-spectrum foundation
or classical Fourier theorem is claimed as a new discovery.

\subsection{The original prime operator and the form domain}

Use the Hilbert space $H=L^2(\R_+^\times,dx/x)$ with inner product
linear in its first variable.  The logarithmic coordinate is $x=e^u$.
Our unitary Fourier transform and original Mellin transform are
\begin{equation}\tag{FMS1}\label{fms:fms:fourier}
 (\mathscr Fg)(t)=\frac1{\sqrt{2\pi}}\int_\R g(e^u)e^{-itu}\,du,
 \qquad \widehat g(t)=\sqrt{2\pi}\,\mathscr Fg(t).
\end{equation}
For a rational prime $p$ put $\ell_p=\log p$ and $r_p=p^{-1/2}$.
Let $(T_aF)(u)=F(u-a)$ on $L^2(\R,du)$.  The norm-convergent series
\begin{equation}\tag{FMS2}\label{fms:fms:bp}
 B_p=(\log p)\sum_{n\ge1}r_p^n(T_{n\ell_p}+T_{-n\ell_p})
\end{equation}
defines a bounded self-adjoint operator, since every translation has
norm one and $\sum_{n\ge1}2r_p^n<\infty$.  The two geometric series give
\begin{equation}\tag{FMS3}\label{fms:fms:multiplier}
 \mathscr FB_p\mathscr F^{-1}=M_{b_p},\qquad
 b_p(t)=(\log p)
 \left(\frac{1-p^{-1}}{1-2p^{-1/2}\cos(t\log p)+p^{-1}}-1\right).
\end{equation}
Here $M_b$ denotes multiplication by $b$.  If
$g^*(x)=\overline{g(x^{-1})}$ and $f=g*g^*$, convolution is always
with respect to $dx/x$.  Direct substitution in the convolution gives
\[
 f(e^a)=\int_\R g(e^u)\overline{g(e^{u-a})}\,du.
\]
The symmetric pair of translations in \eqref{fms:fms:bp} therefore gives
the exact original finite-place functional:
\begin{equation}\tag{FMS4}\label{fms:fms:local}
 \langle B_pg,g\rangle
 =W_p(f)
 =(\log p)\sum_{n\ge1}p^{-n/2}\{f(p^n)+f(p^{-n})\}.
\end{equation}
All terms, including $\log p$, both reciprocal arguments and the
exponent $-n/2$, are retained.

For completeness, this operator is the receiver of the actual
relative-prime energy map in PME6--10, not an assigned multiplier.
Use additive Haar measure $\mu(\Z_p)=1$ on $\Q_p$, and set
\[
 (\mathscr U_p\phi)(y)=\sqrt p\,\phi(y/p),\qquad
 \xi_p=\mathbf1_{\Z_p}\in L^2(\Q_p).
\]
Changing variables proves that $\mathscr U_p$ is unitary, and
$\mathscr U_p^n\xi_p=p^{n/2}\mathbf1_{p^n\Z_p}$ for $n\ge0$.
Hence $\langle\mathscr U_p^n\xi_p,\xi_p\rangle=p^{-|n|/2}$ for
every integer $n$.  For compactly supported $F(u)=g(e^u)$ set
\begin{equation}\tag{FMS5}\label{fms:fms:energy}
 (\mathcal E_pg)(x)=\sum_{n\in\Z}F(x+n\ell_p)\mathscr U_p^n\xi_p,
 \quad 0\le x<\ell_p,\qquad
 \mathcal E_p:H\longrightarrow
 L^2([0,\ell_p);L^2(\Q_p)).
\end{equation}
The sum is initially finite. Expanding its squared norm and
partitioning $\R$ by the intervals $[n\ell_p,(n+1)\ell_p)$ gives
\[
 \|\mathcal E_pg\|^2
 =\sum_{j\in\Z}p^{-|j|/2}
       \int_\R F(u)\overline{F(u-j\ell_p)}\,du.
\]
Cauchy--Schwarz bounds the absolute sum by
$(1+r_p)(1-r_p)^{-1}\|g\|^2$, so $\mathcal E_p$ extends boundedly.
The $j=0$ term is exactly $\|g\|^2$. Polarization now proves
\begin{equation}\tag{FMS6}\label{fms:fms:energyreceiver}
 B_p=(\log p)(\mathcal E_p^*\mathcal E_p-I_H).
\end{equation}
Thus the subtracted identity in \eqref{fms:fms:multiplier} is required by
the original energy receiver.

The denominator of the Poisson kernel in \eqref{fms:fms:multiplier}
varies from $(1-r_p)^2$ to $(1+r_p)^2$.  Consequently
\begin{equation}\tag{FMS7}\label{fms:fms:localrange}
 b_p(\R)=[-\rho_pu_p,u_p],\qquad
 u_p=\frac{2\log p}{\sqrt p-1},\qquad
 \rho_p=\frac{\sqrt p-1}{\sqrt p+1}.
\end{equation}
The extrema occur at $\cos(t\log p)=-1$ and $1$, respectively.

The actual admissible domain is
\begin{equation}\tag{FMS8}\label{fms:fms:domain}
 \mathcal D_0=\{g\in C_c^\infty(\R_+^\times):
                  \widehat g(i/2)=\widehat g(-i/2)=0\}.
\end{equation}
It is dense in $H$.  Indeed the two moment functionals
$F\mapsto\int F(u)e^{u/2}du$ and
$F\mapsto\int F(u)e^{-u/2}du$ are linearly independent on
$C_c^\infty(\R)$: translate a nonnegative nonzero test by a nonzero
real number to obtain an invertible two-by-two moment matrix.
Any linear functional vanishing on their joint kernel therefore is
a linear combination of those two moments.  If an $L^2$ vector
annihilates that kernel, its distribution equals a combination of
$e^{u/2}$ and $e^{-u/2}$.  No nonzero such combination is in $L^2$:
one coefficient produces exponential growth at $+\infty$, and the
other at $-\infty$.  Thus the orthogonal complement of $\mathcal D_0$
is zero.  Algebraic tensor powers of $\mathcal D_0$ are consequently
dense in the corresponding Hilbert tensor powers of $H$.

On $\mathcal D_0$ the complete Weil quadratic form has the convention
\begin{equation}\tag{FMS9}\label{fms:fms:weil}
 Q(g)=-W_{\rm arch}(g*g^*)-\sum_p W_p(g*g^*).
\end{equation}
Each sum on compact tests is finite: only prime powers in a fixed
compact interval can occur in the convolution support.
The Archimedean functional is the original one
\begin{equation}\tag{FMS10}\label{fms:fms:arch}
 W_{\rm arch}(f)
 =(\log(4\pi)+\gamma_E)f(1)+
 \int_1^\infty
 \{f(x)+f(x^{-1})-2x^{-1/2}f(1)\}
       \frac{x^{1/2}}{x-x^{-1}}\frac{dx}{x}.
\end{equation}
For convolution squares the bracket is $O(\log x)$ near $1$, so
the integral is convergent.  Polarization makes all these Hermitian
quadratic forms sesquilinear on $\mathcal D_0$.  No bounded
representing operator for \eqref{fms:fms:weil} will be required.

\subsection{The retained support diagonal and all permutation terms}

Fix distinct primes $I=\{p_1,\ldots,p_s\}$ with $s\ge2$, and let
$\Omega$ be the entire place set.  The actual support faces in this
calculation are the interval
\[
 A\longmapsto(\Omega\setminus I)\cup A,\qquad A\subseteq I.
\]
This map preserves joins and meets and is a lattice isomorphism
onto that interval. Its bottom contains the Archimedean place and
every unselected place.  In particular it is not the external
absent support.  The coefficient receiver used below is the complex
vector space with basis these actual support states:
\begin{equation}\tag{FMS11}\label{fms:fms:diagonal}
 V_I=\C[\mathcal P(I)],\qquad
 h_I=\sum_{A\subseteq I}(-1)^{|A|}\delta_A,\qquad
 \Delta_r:V_I\longrightarrow V_I^{\otimes r},\quad
 \Delta_r\delta_A=\delta_A^{\otimes r}.
\end{equation}
The map $\Delta_r$ is extended linearly.  In particular
$\Delta_rh_I$ is the displayed signed sum of diagonals, and is not
the tensor power $h_I^{\otimes r}$.

Let $q_0$ be the common sesquilinear form
$-W_{\rm arch}-\sum_{p\notin I}W_p$ on convolution pairs from
$\mathcal D_0$, and set
\[
 q_A=q_0-\sum_{p_i\in A}\langle B_{p_i}\,\cdot,\cdot\rangle.
\]
The full support face $A=I$ is the complete form
\eqref{fms:fms:weil}; all other faces retain the very same $q_0$.
Let $\mathcal Q:V_I\to
\{\text{sesquilinear forms on }\mathcal D_0\}$ send
$\delta_A$ to $q_A$.  The precise mixed map is
$\mathcal Q^{\otimes r}\Delta_r$.  Its value on $h_I$ is the form
\begin{equation}\tag{FMS12}\label{fms:fms:mixedform}
 \mathcal J_r=\sum_{A\subseteq I}(-1)^{|A|}q_A^{\otimes r}
 \quad\hbox{on }\mathcal D_0^{\otimes_{\rm alg} r}.
\end{equation}
No product of unbounded operators is present in this definition.

\begin{proposition}[Every tensor summand]\label{fms:fms:subset}
Write $C_0=q_0$ and
$C_i=\langle B_{p_i}\,\cdot,\cdot\rangle$ for $1\le i\le s$.
For a word $w\in\{0,1,\ldots,s\}^r$, let $n(w)$ be its number
of nonzero entries. Then
\begin{equation}\tag{FMS13}\label{fms:fms:allwords}
 \mathcal J_r=
 \sum_{\substack{w\in\{0,\ldots,s\}^r\\
                  \{1,\ldots,s\}\subseteq\{w_1,\ldots,w_r\}}}
       (-1)^{s+n(w)}C_{w_1}\otimes\cdots\otimes C_{w_r}.
\end{equation}
In particular $\mathcal J_r=0$ for $r<s$. At $r=s$ its unique
bounded Hilbert-space receiver is
\begin{equation}\tag{FMS14}\label{fms:fms:mixedoperator}
 \mathbb M_I=\sum_{\sigma\in S_s}
        B_{p_{\sigma(1)}}\otimes\cdots\otimes B_{p_{\sigma(s)}}.
\end{equation}
\end{proposition}
\begin{proof}
In the expansion of $q_A^{\otimes r}$ a fixed word $w$ can occur
exactly when $A$ contains its nonzero labels $T(w)$; its coefficient
there is $(-1)^{n(w)}$. The total coefficient is
\[
 (-1)^{n(w)}\sum_{A\supseteq T(w)}(-1)^{|A|}
 =(-1)^{n(w)+|T(w)|}(1-1)^{s-|T(w)|}.
\]
It vanishes unless $T(w)=I$, in which case it is
$(-1)^{s+n(w)}$.  At $r=s$ each nonzero label must occur exactly
once, there is no zero label, and the coefficient is $1$.
Thus every permutation occurs exactly once. Each factor $B_p$ is
bounded, so their finite tensor sum represents this form. Density
proved after \eqref{fms:fms:domain} gives uniqueness of its bounded
extension. No commutativity between $q_0$ and other forms is needed;
the entries occupy separate tensor factors.
\end{proof}

For $r>s$ formula \eqref{fms:fms:allwords} generally retains the common
form $q_0$. It makes no bounded-extension assertion at those orders.
The finite pulse restrictions and finite traces below exist at every
order without such an assertion.

\subsection{The entire mixed spectrum and its spectral projections}

\begin{theorem}[Exact permanent spectrum]\label{fms:fms:spectrum}
Put $u_i=u_{p_i}$, $\rho_i=\rho_{p_i}$ and
$\rho_I=\max_i\rho_i$.  The bounded self-adjoint operator
$\mathbb M_I$ is unitarily equivalent to multiplication on
$L^2(\R^s)$ by
\begin{equation}\tag{FMS15}\label{fms:fms:permanent}
 m_I(t_1,\ldots,t_s)=
 \perm[b_{p_i}(t_j)]_{i,j=1}^s
 =\sum_{\sigma\in S_s}\prod_{j=1}^s b_{p_{\sigma(j)}}(t_j).
\end{equation}
Its spectrum and norm are
\begin{equation}\tag{FMS16}\label{fms:fms:spectralinterval}
 \sigma(\mathbb M_I)=
 [-s!\rho_I\prod_i u_i,\ s!\prod_i u_i],
 \qquad \|\mathbb M_I\|=s!\prod_i u_i.
\end{equation}
It has no eigenvalues. Its positive and negative spectral subspaces
are both infinite-dimensional, and its kernel is zero.
\end{theorem}
\begin{proof}
The product Fourier transform sends each tensor summand of
\eqref{fms:fms:mixedoperator} to its product multiplier, proving
\eqref{fms:fms:permanent}.  Every factor
$b_{p_i}(t_j)/u_i$ lies in $[-\rho_i,1]$. A nonnegative product
of $s$ such factors is at most $1$.  A negative product has at
least one negative factor; its absolute value is at most
$\rho_I$, since the other factors have absolute value at most $1$.
Termwise summation proves both displayed spectral bounds for $m_I$.

To prove their sharpness with all simultaneous prime frequencies,
write $\mathbb T^s=(\R/2\pi\Z)^s$. The orbit
\[
 t\longmapsto(t\log p_1,\ldots,t\log p_s)\pmod{2\pi}
\]
is dense.  Here is a proof of the exact density needed.
For $n\in\Z^s\setminus\{0\}$ one has
$\sum_i n_i\log p_i\ne0$, since exponentiation and unique prime
factorization forbid $\prod_i p_i^{n_i}=1$. Thus
\[
 \frac1{2T}\int_{-T}^T e^{it\sum_i n_i\log p_i}\,dt
 =\frac{\sin(T\sum_i n_i\log p_i)}
          {T\sum_i n_i\log p_i}\longrightarrow0.
\]
The average of the constant character is $1$. Uniform approximation
by trigonometric polynomials proves that orbit averages of any
continuous function converge to its Haar integral. If a nonempty
open set were missed, choose a nonnegative continuous function
supported inside it, with strictly positive Haar integral.
Its orbit averages would be zero, a contradiction.

The $s$ variables $t_1,\ldots,t_s$ are independent.  Consequently
the corresponding $s$ columns have dense image in
$(\mathbb T^s)^s$.  On this compact connected torus the function
given by the permanent of the Poisson expressions is continuous.
Its maximum is attained by setting every phase equal to $0$.
Its minimum is attained by fixing an index $i_0$ with
$\rho_{i_0}=\rho_I$, setting every phase in row $i_0$ equal to
$\pi$, and every other phase equal to $0$.
Every permutation product is then exactly
$-\rho_I\prod_i u_i$.  These choices attain the two bounds.
The image of a connected space in $\R$ is an interval. Hence the
image on this torus is exactly the interval in
\eqref{fms:fms:spectralinterval}, and the closure of $m_I(\R^s)$ is
that interval.

Every neighborhood of every point in this closed interval has
preimage of positive Lebesgue measure in $\R^s$. Indeed density
first gives a real tuple with its value in that neighborhood,
and continuity gives a nonempty Euclidean open neighborhood of
that tuple with the same property.  Thus the essential range is
the entire interval. Outside it, multiplication by $m_I-z$ has
bounded inverse. At any point of it, restrict a unit vector to a
finite-measure part of the set $|m_I-z|<1/n$ to obtain approximate
eigenvectors. This proves the spectrum and norm assertions.

The denominator in each Poisson expression is strictly positive
on the real axis. Thus $m_I$ is real analytic on $\R^s$ and is
nonconstant, since its range closure contains both endpoints of
a nondegenerate interval. We recall the full measure argument:
the zero set of a nonzero real-analytic function $F$ on $\R^d$
has measure zero. For $d=1$, its zeros are isolated unless it is
identically zero. For the induction step expand in the last
coordinate at zero. At least one coefficient
$\partial_d^kF(x',0)$ is a nonzero analytic function of $x'$;
otherwise analyticity on each vertical line forces $F$ to vanish
everywhere. By induction, the zero set of this one coefficient
has measure zero. Outside it the vertical restriction is not
identically zero and its zero set has one-dimensional measure
zero. Fubini on bounded boxes proves the assertion.
Apply this to $F=m_I-a$ for every real $a$. Therefore an
$L^2$ eigenvector at $a$, which would be supported on
$\{m_I=a\}$, is zero. Nonreal eigenvalues are excluded by
self-adjointness. Both sign regions contain nonempty open sets,
since the range closure has a strictly negative and a strictly
positive value. Their $L^2$ spaces are infinite-dimensional,
for example by using infinitely many disjoint measurable subsets
of a bounded open ball in each region.
\end{proof}

The proof also gives the full spectral measure, not merely its
support. For every Borel set $E\subseteq\R$,
\begin{equation}\tag{FMS17}\label{fms:fms:projection}
 \mathsf E_I(E)=
 (\mathscr F^{\otimes s})^{-1}
       M_{\mathbf1_{\{m_I\in E\}}}\mathscr F^{\otimes s},\qquad
 \langle\mathsf E_I(E)G,G\rangle
 =\int_{\{m_I(t)\in E\}}
            |(\mathscr F^{\otimes s}G)(t)|^2\,dt.
\end{equation}
In particular the sign projections are obtained with
$E=(0,\infty)$ and $E=(-\infty,0)$, and their sum is the identity.
Zero is in the continuous spectrum: $\mathbb M_I$ is injective
with dense range, because its range orthogonal complement is its
kernel, but its range is not closed. If it were closed, its inverse
on that range would be bounded, contradicting the approximate
zero-eigenvectors constructed in the proof.

\subsection{The actual pulse map and its complete compression}

Retain the exact pulse choices of the supplied continuation:
\begin{equation}\tag{FMS18}\label{fms:fms:pulses}
 \begin{gathered}
 M_I=\prod_{i=1}^s p_i,\qquad
 \varepsilon=\frac1{16}\log\frac{M_I+1}{M_I},\\
 w(u)=
 \begin{cases}
 \exp[-1/(1-(u/\varepsilon)^2)],&|u|<\varepsilon,\\
 0,&|u|\ge\varepsilon,
 \end{cases}
 \qquad z=w''-\tfrac14w,\qquad
 N=\int_\R z(u)^2\,du,\\
 a_0=0,\quad a_i=\log p_i,\qquad
 g_c(e^u)=\sum_{j=0}^s c_jz(u-a_j),\quad c\in\C^{s+1}.
 \end{gathered}
\end{equation}
The exponent in $w$ dominates every inverse power of the distance
to its endpoints, so all derivatives vanish at the endpoints.
Thus $w$ and $z$ are smooth. Moreover
$\int zw=-\int|w'|^2-\frac14\int w^2<0$, proving $N>0$.
Two integrations by parts with the zero boundary terms give
\[
 \int_\R z(u)e^{cu}\,du=(c^2-\tfrac14)
               \int_\R w(u)e^{cu}\,du=0,\quad c=\pm\tfrac12.
\]
 Translation multiplies each zero by $e^{ca_j}$.
 Therefore $g_c\in\mathcal D_0$ for every complex coefficient
vector, and every tensor product of these functions retains both
moment conditions in each separate variable.

Put $C(v)=\int z(u)z(u-v)\,du$. It is real and even, vanishes
outside $[-2\varepsilon,2\varepsilon]$, has $C(0)=N$, and satisfies
$|C(v)|\le N$. The pulse centers have pairwise distance greater
than $16\varepsilon$. Indeed, for distinct positive integers
$A,D$ with $D\ge2$,
\begin{equation}\tag{FMS19}\label{fms:fms:integergap}
 |\log(A/D)|\ge\log((D+1)/D).
\end{equation}
For $A>D$ this follows from $A\ge D+1$; for $A<D$ use
$D/(D-1)>(D+1)/D$. Apply this to ratios of distinct primes,
using $D<M_I$; the distances from $0$ are at least $\log2$.
It follows that the map
\begin{equation}\tag{FMS20}\label{fms:fms:embedding}
 \mathcal U:\C^{s+1}\longrightarrow H,\qquad
 (\mathcal Uc)(e^u)=N^{-1/2}\sum_{j=0}^s c_jz(u-a_j),
 \qquad g_c=\sqrt N\,\mathcal Uc
\end{equation}
is an isometry. This changes no original amplitude:
$g_c$ is exactly the unscaled function in \eqref{fms:fms:pulses}.

\begin{proposition}[Lag calculation of every matrix entry]
\label{fms:fms:compression}
For every rational prime $q$,
\begin{equation}\tag{FMS21}\label{fms:fms:primepulse}
 W_q(g_c*g_c^*)=
 \begin{cases}
 2N\beta_i\Re(\overline{c_0}c_i),&q=p_i,\\
 0,&q\notin I,
 \end{cases}
 \qquad \beta_i=\dfrac{\log p_i}{\sqrt{p_i}}.
\end{equation}
Consequently, writing $E_{ab}$ for the matrix taking $e_b$ to
$e_a$ and killing all other basis vectors,
\begin{equation}\tag{FMS22}\label{fms:fms:compressbp}
 \mathcal U^*B_{p_i}\mathcal U=P_i
       =\beta_i(E_{0i}+E_{i0}),\qquad
 \mathcal U^*B_q\mathcal U=0\quad(q\notin I).
\end{equation}
\end{proposition}
\begin{proof}
Expand the convolution with every complex coefficient:
\[
 (g_c*g_c^*)(e^v)=
 \sum_{j,k=0}^s c_j\overline{c_k}\,
                  C(v-a_j+a_k).
\]
Take $v=\log n$ with $n=q^m\ge2$. Terms $j=k$ vanish because
$\log n\ge\log2>2\varepsilon$. If $j=0,k>0$, the argument is
$\log(np_k)>2\varepsilon$. If $j>0,k=0$, the argument is
$\log(n/p_j)$; it is zero only when $n=p_j$, and otherwise
\eqref{fms:fms:integergap} makes its absolute value greater than
$16\varepsilon$. Finally, when $j,k>0$ and $j\ne k$, the
argument is $\log(np_k/p_j)$. It cannot be zero: distinct primes
cannot satisfy $np_k=p_j$ with $n\ge2$. Again
\eqref{fms:fms:integergap}, with denominator $p_j<M_I$, gives a
gap greater than $16\varepsilon$. Hence the only nonzero value
at $v=\log n$ is $Nc_i\overline{c_0}$ when $n=p_i$.
The value at $-\log n$ is its conjugate. Unique factorization
gives $q^m=p_i$ precisely for $q=p_i,m=1$.
Substitution in \eqref{fms:fms:local} proves
\eqref{fms:fms:primepulse}. The norm relation in
\eqref{fms:fms:embedding} and polarization of the Hermitian forms
then prove every entry of \eqref{fms:fms:compressbp}.
\end{proof}

Set $U_s=\mathcal U^{\otimes s}$ and $b=\prod_i\beta_i>0$.
Tensoring the entrywise compression, without suppressing any term,
gives the self-adjoint matrix
\begin{equation}\tag{FMS23}\label{fms:fms:finitemixed}
 \mathfrak M_I=U_s^*\mathbb M_IU_s
       =\sum_{\sigma\in S_s}
          P_{\sigma(1)}\otimes\cdots\otimes P_{\sigma(s)}
 \quad\hbox{on }(\C^{s+1})^{\otimes s}.
\end{equation}

\subsection{Every pulse eigenspace, its inertia, and its leakage}

The orthonormal tensor basis is indexed by words
$w=(w_1,\ldots,w_s)\in\{0,1,\ldots,s\}^s$.
A word is called admissible here when its nonzero labels are
pairwise distinct; this finite-word condition is separate from the
moment condition on test functions. For an admissible word let
$J\subseteq\{1,\ldots,s\}$ be its occupied positions and
$S\subseteq\{1,\ldots,s\}$ its labels. Then $|J|=|S|=k$.
Let $\mathcal H_{J,S}$ be the span of the $k!$ bijective
assignments $J\to S$, with zeros in the other positions.
Define
\begin{equation}\tag{FMS24}\label{fms:fms:blockvectors}
 u_{J,S}=\frac1{\sqrt{k!}}\sum_{w\in\mathcal H_{J,S}}e_w.
\end{equation}
The sum notation means the basis words spanning the indicated block.

\begin{theorem}[Complete pulse spectrum]\label{fms:fms:inertia}
For every complementary pair $(J,S)$ and $(J^c,S^c)$ the vectors
\begin{equation}\tag{FMS25}\label{fms:fms:eigenvectors}
 v_{J,S}^{\pm}=
        \frac{u_{J,S}\pm u_{J^c,S^c}}{\sqrt2}
\end{equation}
are eigenvectors of $\mathfrak M_I$ with respective eigenvalues
\[
 \pm b\sqrt{k!(s-k)!}.
\]
Choose one member of each complementary pair to avoid duplication.
All other directions are zero directions: the zero-sum subspace
in each $\mathcal H_{J,S}$ and the span of words with a repeated
nonzero label. These form its entire kernel.

For $0\le k\le\lfloor s/2\rfloor$ define
\begin{equation}\tag{FMS26}\label{fms:fms:multiplicities}
 \nu_k=
 \begin{cases}
 \binom{s}{k}^2,&k<s/2,\\[1mm]
 \frac12\binom{s}{s/2}^2,&k=s/2\quad(s\text{ even}).
 \end{cases}
\end{equation}
The two eigenvalues of magnitude $b\sqrt{k!(s-k)!}$ each
have multiplicity $\nu_k$. Writing
$n_s=\frac12\binom{2s}{s}$ and
$d_s=(s+1)^s-\binom{2s}{s}$, the exact inertia and
characteristic polynomial are
\begin{equation}\tag{FMS27}\label{fms:fms:inertiapoly}
 (n_+,n_-,n_0)=(n_s,n_s,d_s),\qquad
 \det(\lambda I-\mathfrak M_I)=
 \lambda^{d_s}\prod_{k=0}^{\lfloor s/2\rfloor}
       \bigl(\lambda^2-b^2 k!(s-k)!\bigr)^{\nu_k}.
\end{equation}
Its operator norm and product of all nonzero eigenvalues are
\begin{equation}\tag{FMS28}\label{fms:fms:normpdet}
 \|\mathfrak M_I\|=\sqrt{s!}\,b,\qquad
 \pdet(\mathfrak M_I)=
 (-1)^{n_s}b^{2n_s}
       \prod_{k=0}^{\lfloor s/2\rfloor}
                    [k!(s-k)!]^{\nu_k}.
\end{equation}
\end{theorem}
\begin{proof}
The factor $P_i$ kills $e_j$ unless $j=0$ or $j=i$,
and sends $e_0$ to $\beta_i e_i$ and $e_i$ to $\beta_i e_0$.
For a permutation summand to act nontrivially on a word, the prime
assigned to an occupied position must equal its label.
Two equal nonzero labels therefore force every summand to vanish.
For an admissible word in $\mathcal H_{J,S}$, the assigned prime
is fixed at each occupied position, and the remaining primes may
be assigned bijectively to $J^c$ in exactly $(s-k)!$ ways.
Each such assignment has coefficient $b$ and gives one of the
distinct complementary words. Thus
\begin{equation}\tag{FMS29}\label{fms:fms:blockmap}
 \mathfrak M_I e_w=
 b\sum_{v\in\mathcal H_{J^c,S^c}}e_v,\qquad
 \mathfrak M_Iu_{J,S}
       =b\sqrt{k!(s-k)!}\,u_{J^c,S^c}.
\end{equation}
The matrix between the two word blocks is the constant-entry
matrix $b\,\mathbf1$. Its only nonzero singular value is
$b\sqrt{k!(s-k)!}$; the zero-sum subspaces are killed.
Equivalently, summing the first equality proves the second,
and applying the same argument to the complement gives the
two-by-two matrix
\[
 \begin{array}{c|cc}
 &u_{J,S}&u_{J^c,S^c}\\ \hline
 u_{J,S}&0&b\sqrt{k!(s-k)!}\\
 u_{J^c,S^c}&b\sqrt{k!(s-k)!}&0
 \end{array}.
\]
Its two orthonormal eigenvectors are \eqref{fms:fms:eigenvectors}.
The word blocks and repeated-label words exhaust the tensor
basis, so the listed kernel and eigenspaces are complete.

There are $\binom{s}{k}^2$ pairs $(J,S)$ of size $k$.
For $k<s/2$ each complementary pair has exactly one such member.
At $k=s/2$ complementation has no fixed pair: no subset of a
nonempty position set equals its complement. Hence divide the
count by two in that case. This proves \eqref{fms:fms:multiplicities}.
Counting $s$ chosen elements from two disjoint sets of $s$
elements gives Vandermonde's identity
$\sum_k\binom{s}{k}^2=\binom{2s}{s}$, so
$\sum_{k\le s/2}\nu_k=n_s$.
The total dimension is $(s+1)^s$; this proves the inertia and
characteristic polynomial. Finally
$k!(s-k)!=s!/\binom{s}{k}\le s!$, with equality at $k=0$.
Taking the largest magnitude and multiplying each positive-negative
pair proves \eqref{fms:fms:normpdet}.
\end{proof}

For $s=2$ the inertia is $(3,3,3)$; the nonzero magnitudes
are $\sqrt2\,b$ once per sign and $b$ twice per sign.
For $s=3$ the inertia is $(10,10,44)$; the magnitudes are
$\sqrt6\,b$ once per sign and $\sqrt2\,b$ nine times per sign.
These are evaluations of the general proof, not numerical evidence
substituted for it.

The comparison with the full operator has an exact extra map.
Put $\Pi_s=U_sU_s^*$, the orthogonal projection onto the pulse
tensor space, and define
\begin{equation}\tag{FMS30}\label{fms:fms:leakage}
 L_I=(I-\Pi_s)\mathbb M_IU_s:
       (\C^{s+1})^{\otimes s}\longrightarrow
       (\operatorname{ran}U_s)^\perp.
\end{equation}
Then
\begin{equation}\tag{FMS31}\label{fms:fms:leakidentities}
 \begin{gathered}
 \mathbb M_IU_s=U_s\mathfrak M_I+L_I,\qquad U_s^*L_I=0,\\
 L_I^*L_I=U_s^*\mathbb M_I^2U_s-\mathfrak M_I^2,\qquad
 \|\mathbb M_IU_sv\|^2
       =\|\mathfrak M_Iv\|^2+\|L_Iv\|^2.
 \end{gathered}
\end{equation}
To prove the first line, split the identity into $\Pi_s+(I-\Pi_s)$
and use \eqref{fms:fms:finitemixed}. For the second line multiply
$U_s^*\mathbb M_I(I-\Pi_s)\mathbb M_IU_s$ out, using
$\mathbb M_I=\mathbb M_I^*$ and $U_s^*U_s=I$; orthogonality
gives the norm equality.

The restriction of $L_I$ to every eigenspace of
$\mathfrak M_I$ is injective. Indeed if
$\mathfrak M_Iv=\lambda v$ and $L_Iv=0$, then the first line
of \eqref{fms:fms:leakidentities} gives
$\mathbb M_IU_sv=\lambda U_sv$. The full operator has no
eigenvectors, so $U_sv=0$ and then $v=0$. In particular
\begin{equation}\tag{FMS32}\label{fms:fms:radicalmap}
 L_I:\ker\mathfrak M_I\lhook\joinrel\longrightarrow
       (\operatorname{ran}U_s)^\perp,\qquad
 L_Iv=\mathbb M_IU_sv\quad(v\in\ker\mathfrak M_I).
\end{equation}
Thus every nonzero radical vector in the specified compression
survives under the full operator through this explicit injection.
For an individual normalized compression eigenvector $v$,
$\|(\mathbb M_I-\lambda)U_sv\|=\|L_Iv\|>0$.
The finite and infinite spectral assertions are connected by
\eqref{fms:fms:leakidentities}; neither is substituted for the other.

\subsection{All higher tensor orders on the same tests}

Here the common term on the pulse space is exactly the
Archimedean term, since \eqref{fms:fms:primepulse} proves that every
unselected finite prime is zero. To retain the exact coefficient
receiver, put
\[
 K(t)=\frac{e^{-t/2}}{1-e^{-2t}},\qquad
 c_{\rm arch}=\log(4\pi)+\gamma_E,
\]
and define the convergent integrals
\begin{equation}\tag{FMS33}\label{fms:fms:archcoeff}
 \begin{split}
 A_\varepsilon
 &=c_{\rm arch}N+
       2\int_0^\infty[C(t)-Ne^{-t/2}]K(t)\,dt,\\
 J_\varepsilon(d)
 &=\int_{d-2\varepsilon}^{d+2\varepsilon}
                    C(t-d)K(t)\,dt\quad(d>16\varepsilon).
 \end{split}
\end{equation}
The integrand of the first line is bounded near zero, since
$C(t)-Ne^{-t/2}=O(t)$ and $K(t)=O(t^{-1})$.
The original negative subtraction is exponentially integrable
outside the support of $C$.
Expansion of \eqref{fms:fms:arch}, using the convolution in the proof
of Proposition~\ref{fms:fms:compression}, gives
\begin{equation}\tag{FMS34}\label{fms:fms:facecoeff}
 \begin{gathered}
 W_{\rm arch}(g_c*g_c^*)=
 A_\varepsilon\sum_{j=0}^s|c_j|^2+
       2\sum_{0\le j<k\le s}
            \Re(\overline{c_j}c_k)J_\varepsilon(|a_j-a_k|),\\
 \mathsf W_A=-A_\varepsilon I
       -[J_\varepsilon(|a_j-a_k|)]_{j\ne k}
       -N\sum_{p_i\in A}\beta_i(E_{0i}+E_{i0}),\\
 Q_A(g_c)=c^*\mathsf W_Ac.
 \end{gathered}
\end{equation}
The off-diagonal matrix in the second line has zero diagonal.
For each pair of different centers, one of the two shifted
correlations is supported on the positive half-axis and the other
on the negative half-axis. The two reciprocal values in
\eqref{fms:fms:arch} therefore yield precisely the factor
$2\Re(\overline{c_j}c_k)$ in the first line, with no assumption
that $c$ is real. The same formula follows when two positive
lag intervals overlap, since it is a linear integral of their sum.

The companion finite-pulse theorem \cite{FinitePrimePulse}, FP37 with its complete
proof FP38--43, with exactly
\eqref{fms:fms:pulses} and \eqref{fms:fms:archcoeff}, proves the quantitative
matrix inequality
\begin{equation}\tag{FMS35}\label{fms:fms:finitebound}
 \mathsf W_A\succeq\mu_*NI_{s+1}\quad(A\subseteq I),\qquad
 \mu_*=\frac{6600077509}{144027072000}>\frac1{22}.
\end{equation}
This is a proved theorem about the entire $(s+1)$-dimensional
complex pulse coefficient space and every finite support face.
It is the only positivity input below. It makes no assertion on
other compact tests or on a bounded global Weil operator.
The companion proof evaluates the full subtraction in
\eqref{fms:fms:archcoeff}; it is not an assumption that that integral
has the desired sign.

For the original unscaled $g_c$ put
\begin{equation}\tag{FMS36}\label{fms:fms:cornerparameters}
 A_c=-W_{\rm arch}(g_c*g_c^*),\qquad
 \alpha_i=2N\beta_i\Re(\overline{c_0}c_i),\qquad
 Q_A(g_c)=A_c-\sum_{p_i\in A}\alpha_i.
\end{equation}
The symbol $\alpha_i$ is the quantity denoted $a_i$ in the
supplied continuation's Section 8; it is not the logarithmic
pulse center $a_i$ in \eqref{fms:fms:pulses}. This distinction changes
no parameter or formula.
For integers $r\ge0$ define the literal support observation
\begin{equation}\tag{FMS37}\label{fms:fms:jr}
 J_r(c)=\sum_{A\subseteq I}(-1)^{|A|}
                   \left(A_c-\sum_{p_i\in A}\alpha_i\right)^r.
\end{equation}
At $r=0$ each power is defined to be $1$, as for a zeroth tensor
power; the signed sum is zero since $s>0$.

\begin{theorem}[All orders, signs, bounds, and exact rate]
\label{fms:fms:higher}
For every $c\in\C^{s+1}$, $J_r(c)=0$ for $r<s$.
For $r\ge s$,
\begin{equation}\tag{FMS38}\label{fms:fms:integral}
 J_r(c)=\frac{r!}{(r-s)!}\left(\prod_{i=1}^s\alpha_i\right)
 \int_{[0,1]^s}
          \left(A_c-\sum_{i=1}^s t_i\alpha_i\right)^{r-s}\,dt.
\end{equation}
For nonzero $c$, set
\begin{equation}\tag{FMS39}\label{fms:fms:minmax}
 m_c=A_c-\sum_i(\alpha_i)_+,\qquad
 M_c=A_c+\sum_i(-\alpha_i)_+.
\end{equation}
These are the smallest and largest actual support-face values,
and $m_c\ge\mu_*N\|c\|_2^2>0$.
For every $r\ge s$ one has
\begin{equation}\tag{FMS40}\label{fms:fms:sign}
 \sgn J_r(c)=
   \sgn\prod_{i=1}^s\Re(\overline{c_0}c_i),
\end{equation}
with $J_r(c)=0$ exactly when at least one factor on the right
vanishes. The complete quantitative bounds are
\begin{equation}\tag{FMS41}\label{fms:fms:bounds}
 \frac{r!}{(r-s)!}\left|\prod_i\alpha_i\right|m_c^{\,r-s}
 \le |J_r(c)|\le
 \frac{r!}{(r-s)!}\left|\prod_i\alpha_i\right|M_c^{\,r-s}.
\end{equation}
If every $\alpha_i$ is nonzero, then
\begin{equation}\tag{FMS42}\label{fms:fms:rate}
 \lim_{r\to\infty}|J_r(c)|^{1/r}=M_c,\qquad
 \frac{|J_r(c)|}{M_c^r}\longrightarrow1.
\end{equation}
\end{theorem}
\begin{proof}
For a differentiable function $F$ and a real $\alpha$, including
negative $\alpha$, the fundamental theorem of calculus gives
\[
 F(x)-F(x-\alpha)=
       \alpha\int_0^1F'(x-t\alpha)\,dt.
\]
Apply this identity successively with
$\alpha=\alpha_1,\ldots,\alpha_s$ to $F(x)=x^r$ at $x=A_c$.
The left side expands to exactly \eqref{fms:fms:jr}, with its original
sign $(-1)^{|A|}$. The $s$th derivative is zero for $r<s$, and
is $r!(r-s)!^{-1}x^{r-s}$ otherwise. This proves both assertions
through \eqref{fms:fms:integral}, also when some $\alpha_i=0$.

The affine function on the cube in \eqref{fms:fms:integral} is
minimized by $t_i=1$ for $\alpha_i>0$ and $t_i=0$ for
$\alpha_i<0$; its maximum uses the reverse choices.
This proves \eqref{fms:fms:minmax}, and these choices are cube
vertices, hence actual face values. The finite theorem
\eqref{fms:fms:finitebound} gives the stated positive lower bound.
The integral in \eqref{fms:fms:integral} is therefore strictly positive
for every $r\ge s$. Since $N$ and every $\beta_i$ are strictly
positive, the sign of its prefactor is exactly
\eqref{fms:fms:sign}. Bounding the integrand by its minimum and
maximum, on a cube of volume one, proves \eqref{fms:fms:bounds}.
At $c=0$ all $J_r$ with $r>0$ are zero, so no separate positivity
claim is required.

When every $\alpha_i\ne0$, the unique maximizing face is
$A_{\max}=\{p_i:\alpha_i<0\}$, and its sign in
\eqref{fms:fms:jr} is
$\epsilon_c=(-1)^{|A_{\max}|}=\sgn\prod_i\alpha_i$.
All other corner values are strictly between $0$ and $M_c$.
Dividing \eqref{fms:fms:jr} by $M_c^r$ proves convergence to
$\epsilon_c$ and hence both limits. More explicitly, if
$d_c=\min_i|\alpha_i|>0$, every other corner is at most
$M_c-d_c$, and all corners are positive. Thus the exact finite
sum supplies the error bound
\begin{equation}\tag{FMS43}\label{fms:fms:rateerror}
 |J_r(c)-\epsilon_c M_c^r|
       \le(2^s-1)(M_c-d_c)^r.
\end{equation}
\end{proof}

For literal tests realizing all three possibilities, take $c_0=1$.
The choice $c_i=1$ for every $i$ gives $J_r(c)>0$ for every
$r\ge s$. The choice $c_1=-1$ and $c_i=1$ for $i>1$ gives
$J_r(c)<0$ at every such order. The choice $c_1=i$ and
$c_j=1$ for $j>1$ gives $J_r(c)=0$ at every order, while
$g_c\ne0$ and every face form remains strictly positive by
\eqref{fms:fms:finitebound}. These are the original amplitude choices
in \eqref{fms:fms:pulses}, including the complex phase in the third test.

These formulas retain all polynomial coefficients. An additional
way to see them simultaneously is the exact exponential generating
identity and its coefficient formula:
\begin{equation}\tag{FMS44}\label{fms:fms:generating}
 \begin{split}
 \sum_{r=0}^\infty\frac{J_r(c)}{r!}z^r
     &=e^{A_cz}\prod_{i=1}^s(1-e^{-\alpha_i z}),\\
 J_r(c)&=
 \sum_{\substack{n_0\ge0,\ n_i\ge1\\n_0+n_1+\cdots+n_s=r}}
 \frac{r!}{n_0!n_1!\cdots n_s!}
 A_c^{n_0}\prod_{i=1}^s
                   (-1)^{n_i+1}\alpha_i^{n_i}.
 \end{split}
\end{equation}
The first identity follows by summing the finite subset exponential
and factoring; expanding the $s+1$ entire power series proves the
second. In particular
$J_s(c)=s!\prod_i\alpha_i$ and
$J_{s+1}(c)=(s+1)!\prod_i\alpha_i
             (A_c-\frac12\sum_i\alpha_i)$.
The first formula agrees exactly with the compression, including
its original mass:
\begin{equation}\tag{FMS45}\label{fms:fms:nfactor}
 J_s(c)=N^s
 \langle\mathfrak M_Ic^{\otimes s},c^{\otimes s}\rangle
 =s!(2N)^s\left(\prod_i\beta_i\right)
                  \prod_i\Re(\overline{c_0}c_i).
\end{equation}
Indeed the matrix of $q_A$ on $\mathcal U\C^{s+1}$ is
$N^{-1}\mathsf W_A$. Proposition~\ref{fms:fms:subset} at order $s$
and \eqref{fms:fms:finitemixed} give the first equality. The product
of the individual star-matrix values gives the second.

Finally all orders have a finite trace receiver that preserves the
support diagonal.  Let $\Pi_A$ be projection onto $\C\delta_A$
in $V_I$, let $D_{h_I}\delta_A=(-1)^{|A|}\delta_A$, and for a
finite vector $v$ put $R_vx=\langle x,v\rangle v$. Then
\begin{equation}\tag{FMS46}\label{fms:fms:trace}
 \begin{gathered}
 \mathsf O_r=\sum_{A\subseteq I}\Pi_A\otimes\mathsf W_A^{\otimes r},
 \\
 \Tr\left[(D_{h_I}\otimes R_{c^{\otimes r}})\mathsf O_r\right]
       =J_r(c).
 \end{gathered}
\end{equation}
For one support block, the rank-one trace is
$(c^*\mathsf W_Ac)^r$, because
$\Tr(TR_v)=\langle Tv,v\rangle$; an orthonormal basis expansion
proves this identity directly. Summing the support blocks proves
\eqref{fms:fms:trace}. Each $\mathsf W_A$ is positive by the proved
finite bound, so $\mathsf O_r$ is positive. The actual inserted
$D_{h_I}$ has both signs and is not a positive functional.
Consequently both signs in \eqref{fms:fms:sign} are compatible with
the positive face matrices through this exact trace map.
This finite receiver does not take an unspecified infinite trace.

\paragraph{Precise scope of the result.}
The full Hilbert-space assertion is the evaluated bounded operator
\eqref{fms:fms:mixedoperator}, with every spectral projection in
\eqref{fms:fms:projection}. The finite tensor eigenspaces are its exact
compression through \eqref{fms:fms:embedding}, with the complementary
map \eqref{fms:fms:leakage}. The all-order trace keeps the same support
diagonal and the exact common Archimedean number. These maps prove
how the different sign calculations are related. They establish no
global Weil positivity theorem, no zero off the critical line, and
no new purity statement.


\begin{figure}[p]
\centering
\begin{tikzcd}[column sep=large,row sep=large]
 (\C^{s+1})^{\otimes s}
 \arrow[r,"U_s",hook]\arrow[d,"\mathfrak M_I"']
 &H^{\otimes s}\arrow[d,"\mathbb M_I"]\\
 (\C^{s+1})^{\otimes s}\arrow[r,"U_s",hook]
 &H^{\otimes s}\arrow[l,bend left=20,"U_s^*"']
\end{tikzcd}

\medskip
The square has the exact defect
\[
 \mathbb M_IU_s-U_s\mathfrak M_I=L_I,
 \qquad U_s^*L_I=0,
 \qquad L_I^*L_I=U_s^*\mathbb M_I^2U_s-\mathfrak M_I^2.
\]

\begin{minipage}{0.92\linewidth}
\textbf{The full operator.}
Its spectrum is the entire interval
\[
 [-s!\rho_I\prod_i u_i,\ s!\prod_i u_i],
 \qquad u_i=\frac{2\log p_i}{\sqrt{p_i}-1},\quad
 \rho_I=\max_i\frac{\sqrt{p_i}-1}{\sqrt{p_i}+1}.
\]
Both signs have infinite-dimensional spectral subspaces. There are
no eigenvectors, including at zero.

\medskip
\textbf{The complete finite compression for two primes.}
Put $b=(\log p_1)(\log p_2)/\sqrt{p_1p_2}$. Every one of
its nine dimensions is represented in this table:
\[
\begin{array}{c|rrrrr}
 \text{eigenvalue}&-\sqrt2b&-b&0&b&\sqrt2b\\\hline
 \text{multiplicity}&1&2&3&2&1
\end{array}
\]
The three-dimensional zero eigenspace maps injectively through
$L_I$ into the orthogonal complement of the pulse space. This is
the calculated fate of its apparent zero directions under the
original full operator.

\medskip
\textbf{The Weil face matrices on the same tests.}
Every face retains its original Archimedean integral and obeys
$\mathsf W_A\succeq\mu_*N I$, where
$\mu_*=6600077509/144027072000>1/22$.
Their mixed tensor trace inserts $D_{h_I}$, whose diagonal entries
are $(-1)^{|A|}$. The sign in the mixed trace comes from this
specified insertion; the face matrices remain positive.
\end{minipage}
\caption{The exact maps from the finite tests to the full mixed-prime
operator, including the square's nonzero defect. FMS14--17 prove
the full spectrum, FMS20--32 give the compression and complementary
map, and FP37--43 and FMS46 prove the positive face bound and
signed trace. All original pulse amplitudes are recovered by
$g_c=\sqrt N\,\mathcal Uc$. The local analytic distributions
are those of Connes \cite{Connes2026}, with the programme's
prime-energy receiver in the complete source \cite{Edition033}.}
\end{figure}
